\documentclass[palette=none,mode=dark]{impaTeXPoster}

\usepackage[utf8]{inputenc}
\usepackage[T1]{fontenc}
\usepackage{lipsum}

\title{Poster Teste}
\author{Bruno Pereira, Maria Silva, João Souza}
\authoremail{bruno@impa.br, maria@impa.br, joao@impa.br}
\leftposterlogo{images/impa.png}
\posterlogo{images/logo.png}

\begin{document}
\maketitle

\begin{columns}

% ==========================================
% LEFT COLUMN
% ==========================================
\column{0.5}

\begin{postercard}[black,black,white,orange]{Introdução}
Este é um exemplo de seção usando a classe ImpaTeXPoster.

\lipsum[1]
\end{postercard}

\begin{postercard}[black,black,white,orange]{Objetivo}
Mostrar como usar diferentes cores no ambiente.

\lipsum[2]
\end{postercard}

% ==========================================
% RIGHT COLUMN
% ==========================================
\column{0.5}

\begin{postercard}[black,black,white,orange]{Resultados}
\lipsum[3]
\end{postercard}

% primeiro parâmetro é a cor do Fundo do título
% segundo parâmetro é a cor do do testo
% terceiro parâmetro é cor do texto
% quarto parâmetro é a borda do postercard

\begin{postercard}[black,black,white,orange]{Resultados}
\lipsum[4]
\end{postercard}

\begin{postercard}[black,black,white,orange]{Referências}

\begin{thebibliography}{99}

\bibitem{popper}
POPPER, Karl.
\textit{A lógica da pesquisa científica}.
São Paulo: Cultrix, 2002.

\bibitem{chalmers}
CHALMERS, A. F.
\textit{O que é ciência, afinal?}
São Paulo: Brasiliense, 1999.

\end{thebibliography}

\end{postercard}

\end{columns}

\posterfooter{https://github.com/impaTeX/modelo-cartaz/}{%
    \begin{thebibliography}{99}
    \bibitem{popper}
    POPPER, Karl.
    \textit{A lógica da pesquisa científica}.
    São Paulo: Cultrix, 2002.

    \bibitem{chalmers}
    CHALMERS, A. F.
    \textit{O que é ciência, afinal?}
    São Paulo: Brasiliense, 1999.
    \end{thebibliography}
}

\end{document}